\documentclass[10pt,a4paper]{amsart}
\usepackage[margin=19mm]{geometry}
\usepackage{amsmath,amssymb,mathtools}
\usepackage[scr=rsfs,cal=euler]{mathalfa}
\usepackage{xfrac}
\usepackage[unicode,hidelinks,bookmarks=false]{hyperref}
\newcommand{\ie}{i.e.\ }
\newcommand{\cf}{cf.\ }
\newcommand{\resp}{respectively\ }
\newcommand{\Subsection}{\subsection}
\providecommand{\nobreakdash}{\nobreak\hbox{-}}
\setlength{\emergencystretch}{2em}
\multlinegap0pt
\usepackage{amssymb}
\newtheorem{theorem}{Theorem}
\title{Advanced linear algebra}
\author{Teo Banica}
\date{Source: arXiv:2506.18666v2 (CC BY 4.0)}
\begin{document}
\maketitle

\noindent\emph{Source and licence.} Advanced linear algebra. Version arXiv:2506.18666v2 (CC BY 4.0), distributed by arXiv under the Creative Commons Attribution 4.0 International licence, http://creativecommons.org/licenses/by/4.0/, as recorded in the arXiv licence metadata for this exact version and shown on the abstract page https://arxiv.org/abs/2506.18666v2. Full licence notice: \texttt{licenses/arxiv-2506.18666-CC-BY-4.0.txt}.\par\smallskip

\noindent\textit{Editorial scope.} Counted unit: the article's theorem theorem (source0.tex lines 6011--6015). The article's complete original proof of it is delivered with it (source0.tex lines 6017--6153, one \texttt{proof} environment of 137 lines). No mathematics is written, repaired or re-derived: the statement and the proof are the article's own lines, in the article's own notation, and no retained line is substituted or re-flowed.  No further source-local context is needed: every cross-reference of every retained line resolves inside the two delivered spans.  Nothing was omitted from any retained span, so the package carries no omission record.  No retained line contains a citation, so the wrapper carries no bibliography and no bibliography entry.  No retained line contains a figure environment, an image-inclusion command or drawing code, so no image is delivered and the package declares no asset dependency.  Editorial formatting only, in addition: an AMS \texttt{amsart} preamble that restates the article's own declarations and macros used by the retained lines, namely the environments \texttt{theorem} on separate counters, together with the standard packages those lines need.  The retained environments print in this document as Theorem 1 in order of appearance, whereas the article prints the same items with the numbering of its own full document.  The complete native source of the article as distributed in the arXiv e-print for this exact version is bundled unchanged as \texttt{sources/180353/source0.tex}.
\begin{theorem}
We can apply $f(x)=x^2$ to any matrix $A\in M_N(\mathbb C)$, 
$$A\to A^2$$
and if $\lambda_1,\ldots,\lambda_N$ are the eigenvalues of $A$, then $\lambda_1^2,\ldots,\lambda_N^2$ are the eigenvalues of $A^2$.
\end{theorem}

\begin{proof}
This does not look difficult to establish, but let us have a detailed discussion about this, which will be quite instructive, by following Method 6.7: 

\medskip

(1) Starting with bare hands, let us see if we can solve the problem at $N=2$, via a direct computation, without using any trick. So, consider a $2\times2$ matrix, as follows:
$$A=\begin{pmatrix}
a&b\\ 
c&d
\end{pmatrix}$$

The square of this matrix is then given by the following formula:
$$A^2=\begin{pmatrix}
a&b\\ 
c&d
\end{pmatrix}\begin{pmatrix}
a&b\\ 
c&d
\end{pmatrix}
=\begin{pmatrix}
a^2+bc&ab+bd\\ 
ac+cd&bc+d^2
\end{pmatrix}$$

Now let us compare the eigenvalues of the matrices $A,A^2$. Those of the initial matrix $A$, say $r,s\in\mathbb C$, are subject to the following two equations:
$$r+s=Tr(A)=a+d$$
$$rs=\det(A)=ad-bc$$

As for the eigenvalues of the squared matrix $A^2$, say $R,S\in\mathbb C$, these are subject to some similar equations, which are as follows:
$$R+S=Tr(A^2)=a^2+d^2+2bc$$
$$RS=\det(A^2)=\det(A)^2$$

The second equality suggests that we should have $R=r^2,S=s^2$, so let us prove now that it is indeed so. For this purpose, we just have to check that the numbers $R=r^2,S=s^2$ sum up to the quantity computed above, and this is done as follows:
\begin{eqnarray*}
R+S
&=&r^2+s^2\\
&=&(r+s)^2-2rs\\
&=&(a+d)^2-2(ad-bc)\\
&=&a^2+d^2+2bc
\end{eqnarray*}

Summarizing, result proved with bare hands at $N=2$. However, when thinking a bit, such methods will become quite complicated at $N\geq3$, so we will stop here, with this.

\medskip

(2) Still with bare hands, but allowing us some tricks, namely the use of square roots of complex numbers, here is how the result can be established, at any $N\in\mathbb N$. Consider our matrix $A\in M_N(\mathbb C)$, and let us factorize its characteristic polynomial, as follows:
$$\det(A-x)=\prod_i(\lambda_i-x)$$

We have then the following computation, for the characteristic polynomial of $A^2$:
\begin{eqnarray*}
\det(A^2-x)
&=&\det((A-\sqrt{x})(A+\sqrt{x}))\\
&=&\det(A-\sqrt{x})\det(A+\sqrt{x})\\
&=&\prod_i(\lambda_i-\sqrt{x})\prod_i(\lambda_i+\sqrt{x})\\
&=&\prod_i(\lambda_i-\sqrt{x})(\lambda_i+\sqrt{x})\\
&=&\prod_i(\lambda_i^2-x)
\end{eqnarray*}

Thus, claimed proved, eventually, with bare hands, or almost.

\medskip

(3) Getting now to more advanced tricks, bringing heavy simplifications, we can use here density arguments. Indeed, assume first that our matrix is diagonalizable:
$$A=P\begin{pmatrix}
\lambda_1\\
&\ddots\\
&&\lambda_N
\end{pmatrix}
P^{-1}$$

In this case, we have the following computation, for the squared matrix:
\begin{eqnarray*}
A^2
&=&P\begin{pmatrix}
\lambda_1\\
&\ddots\\
&&\lambda_N
\end{pmatrix}
P^{-1}\cdot P\begin{pmatrix}
\lambda_1\\
&\ddots\\
&&\lambda_N
\end{pmatrix}
P^{-1}\\
&=&P\begin{pmatrix}
\lambda_1\\
&\ddots\\
&&\lambda_N
\end{pmatrix}^2
P^{-1}\\
&=&P\begin{pmatrix}
\lambda_1^2\\
&\ddots\\
&&\lambda_N^2
\end{pmatrix}
P^{-1}
\end{eqnarray*}

Thus, claim proved in this case, and the general case follows now by using the fact, that we know well from chapter 4, that the diagonalizable matrices are dense.

\medskip

(4) Finally, for our discussion to be complete, let us see as well what we get, by using a nuclear bomb, that is, the Jordan form. So, let us write the matrix to be squared in Jordan form, as in chapter 5, as follows, with $P$ denoting the passage matrix:
$$A=P\begin{pmatrix}
J_1\\
&\ddots\\
&&J_k
\end{pmatrix}P^{-1}$$

The square of this matrix is then given by the following formula:
$$e^A=P\begin{pmatrix}
J_1^2\\
&\ddots\\
&&J_k^2
\end{pmatrix}P^{-1}$$

Thus, it is enough to know how to square the Jordan blocks. So, consider a Jordan block, as follows, with our usual convention that blank spaces stand for $0$ entries:
$$J=\begin{pmatrix}
\lambda&1\\
&\lambda&1\\
&&\ddots&\ddots\\
&&&\lambda&1\\
&&&&\lambda
\end{pmatrix}$$

The square of this Jordan block is then given by the following formula:
$$J=\begin{pmatrix}
\lambda^2&2\lambda&1\\
&\lambda^2&2\lambda&1\\
&&\ddots&\ddots&\ddots\\
&&&\lambda^2&2\lambda&1\\
&&&&\lambda^2&2\lambda\\
&&&&&\lambda^2
\end{pmatrix}$$

Thus, we have the square $\lambda^2$ of the eigenvalue $\lambda$ on the diagonal, and by putting everything together, we are again led to the conclusion in the statement.
\end{proof}

\end{document}
